\section{Scope and Imported Mathematical Spine}\label{sec:scope-and-spine}

\subsection{Downstream sector scope and calculation spine}\label{subsec:sector-scope}

\paragraph{Setup.}
This note instantiates the foundational calculus of \mainref{v41.02}~\cite{poyau-aod-v4102} on the coord-0 elemental support ladder $Z = 1 \dots 118$ (together with the unretained open cutoff at $Z = 119 \dots 126$ on $\mathsf{C}_{22}$). No new axioms or primitive relations are introduced. Every calculation begins at the first finite anchor $a_0 = \mathtt{00}_8$ (\mainref{(3)}) with local origin $0\vdash$, declares a tesseract-admissible boundary and integer comparison window
\begin{equation}
\mathbf{a} \vdash (F, R, \partial_T R, \omega), \qquad B = (\partial_T F, \omega),
\label{eq:matter1-scope-binder}
\end{equation}
and executes the active $\AO$ calculation spine (\mainref{\S5.13, Fig.~11}; Manual~\S1.3):
\begin{equation}
C_f[\Sigma]_{\AO}, \qquad \AOh(C) = \langle \Pi_\alpha(C), H_\mu(C), \Pi_\Omega(C)\rangle, \qquad \SADARop_B(C) = p^D_{C,\AO}(B) A_{C,\AO}(B),
\label{eq:matter1-calculation-spine}
\end{equation}
with boundary accumulation
\begin{equation}
\Phi^D_{\partial_T R} = \sum_{e \in \partial_T R} \SADARop_B(C_e).
\label{eq:matter1-boundary-accumulation}
\end{equation}
Tables in this note store evaluated integer and rational values produced by \eqnref{eq:matter1-scope-binder}--\eqnref{eq:matter1-boundary-accumulation}; tables are document containers, not the calculus.

\paragraph{Derivation.}
By \mainref{(1)--(5)}, Null sameness ($\Null \equiv \infty,\; \Null \to \Null \to \Null \equiv \Null$) is schematic idempotency; once a finite distinction commits an event on $Q_4 = \{0,1\}^4$ (\mainref{(8)}), its identity, lineage, and branch/hinge/counterbranch order $M_{\varepsilon,m} = \varepsilon\,0_H(0_D)^m(-\varepsilon)$ remain in the trace. By \mainref{(12)--(16)}, a certified closure $\operatorname{Cert}(C)$ requires at least two committed phase-bearing traversals ($m(C) \ge 2$) together with its retained-start return witness $\mathcal{W}_C$, contributing one closure-phase cycle $\Delta\Phi(C) = 1_C$.

\paragraph{Specialization.}
In the $D0$ exact internal fixture layer of this note, all support cardinalities $\Pi$, reflection durations $RD$, channel pressures $P^D$, closure compatibilities $C_{\mathrm{close}}$, rung residuals $S_{n+1}$, and sheddic surpluses $X_{\mathrm{shedding}}$ are exact integers in $\mathbb{Z}$. Measured-sector element symbols ($\mathrm{H} \dots \mathrm{Og}$ for $Z = 1 \dots 118$) are application-layer carrier labels attached only after the coord-0 construction is admitted (\mainref{\S5.17.10}).

\paragraph{Calibration.}
Each elemental rung binds one integer coord-0 index $Z \in \{1, \dots, 118\}$ and one reference nucleon count $A \in \mathbb{Z}_{\ge 1}$ in the count basis $\mathcal{B}_{\mathrm{curl}} = (3{:}3, 3{:}4, 1{:}2)$ (\mainref{App.~B.4.1}), with collision/channel capacity $C_3 = 3$ (\mainref{(77)}, Manual~\S5.4). External mass or spectral target comparisons remain quarantined in $\Qext(\nu_{x,\AO}; \mathcal{T}_x)$ (Manual~\S1.8, App.~B).

\paragraph{Falsification.}
Importing an undeclared external constant into an $\AO$-internal calculation, collapsing Null $\Null$, fibre zero $0$, vertex $0000$, anchor $\mathtt{00}_8$, and scope $B_0$ into a single type, or using a table label without a bound closure trace $HT(K;B)$ contradicts \eqnref{eq:matter1-scope-binder}--\eqnref{eq:matter1-boundary-accumulation}.

\aodprovenance{\mainref{(1)--(5)}, \mainref{(8)}, \mainref{(11)--(16)}, \mainref{(40)}, \mainref{(44)}, \mainref{(92)}.}
\aodliteraturenote{Foundational protocol and calculus: Poyau~\cite{poyau-af,poyau-afc,poyau-aod-v4102}.}

\subsection{Curling-curl support and branch/hinge/branch reflection duration}\label{subsec:curling-curl-rd}

\paragraph{Setup.}
A support expression $\Sigma = (s_1, \dots, s_m; L)$ (\mainref{Def.~8, (36)}) with clauses $s_i = p_i{:}q_i$ and optional shared enclosure qualifier $L$ proposes a curling-curl support specification $K_{\mathrm{prop}}$. Admission on $B = (\partial_T F, \omega)$ is the traced calculation (\mainref{Def.~9, (37)})
\begin{equation}
K_{\mathrm{prop}} \xrightarrow{\,HT(K;B)\,} K_{\mathrm{adm}}(B),
\label{eq:matter1-traced-admission}
\end{equation}
binding the complete motif record $K_{\mathrm{adm}}(B,i;\mathrm{cert})$ of \mainref{(33)}. For each curling-curl clause $p{:}q$, import the finite edge-shell walk support model of \mainref{\S3.5, (35)} and \mainref{App.~H.5, (168)}:
\begin{equation}
\Pi(p{:}q) = q^p + q, \qquad p, q \in \mathbb{Z}_{\ge 1}.
\label{eq:matter1-pi-support}
\end{equation}
Here $q^p$ counts memoryless $q$-direction walk branches and $+q$ is the retained Hamming-1 edge-shell term on $Q_4$.

\paragraph{Derivation.}
For a certified cycle $C$ on $B$, the Alpha--Omega hinge kernel (\mainref{Def.~14, (44)}) has branch/hinge/branch reflection duration (\mainref{\S3.3.2, (31)}, \mainref{\S3.9, (49)}, and \mainref{App.~H.5, (169)}):
\begin{equation}
\AOh(C) = \langle \Pi_\alpha(C), H_\mu(C), \Pi_\Omega(C)\rangle,
\qquad
RD_{\AO}(C;B) = |\Pi_\alpha(C)| + H_\mu(C) + |\Pi_\Omega(C)|.
\label{eq:matter1-rd-general}
\end{equation}
Under symmetric branch supports ($|\Pi_\alpha(C)| = |\Pi_\Omega(C)| = \Pi(p{:}q)$), \eqnref{eq:matter1-rd-general} gives $RD_{\AO}(p{:}q;B) = 2\Pi(p{:}q) + H_\mu(B)$ (\mainref{(32)}). Under the single-hinge reduction $H_\mu(B) = 1$ (\mainref{(50), (170)}),
\begin{equation}
RD^{(1)}_{\AO}(p{:}q) = 2\Pi(p{:}q) + 1 = 2(q^p + q) + 1.
\label{eq:matter1-rd-single-hinge}
\end{equation}

\paragraph{Specialization.}
Evaluating \eqnref{eq:matter1-pi-support} and \eqnref{eq:matter1-rd-single-hinge} on the canonical curling-curl families of \mainref{App.~B.2, Table~1} yields the exact integer primitives in \tabref{tab:matter1-curling-curl-primitives}. Note the strict typing separation of \mainref{(34)} and \mainref{App.~B.1.10}: $\mathrm{Dimonon} = 2\,\mathrm{Monon} \ne \mathrm{Dion}$, and $\mathrm{tetrion}$ (bare four-cycle), $\mathrm{tetron}$ ($E_4$ four-seat inner enclosure), $\mathrm{tetratrion}$ ($4\times\mathrm{trion} = 3{:}4$), and $\mathrm{tetrad}$ (seated tetron closure) are distinct types.

\begin{table}[H]
\centering
\footnotesize
\setlength{\tabcolsep}{4.5pt}
\caption{\textbf{Imported curling-curl support primitives and single-hinge reflection durations (\mainref{\S3.5}, \mainref{App.~B.2}, \mainref{App.~H.5}).}}
\label{tab:matter1-curling-curl-primitives}
\begin{tabular}{@{}llllrrl@{}}
\toprule
Form $p{:}q$ & Family name & Stem & Meaning & $\Pi(p{:}q)=q^p+q$ & $RD^{(1)}_{\AO}=2\Pi+1$ & Default hexon instance ($\mathsf{C}_6$)\\
\midrule
$1{:}2$ & Dimonon & Dimon & $2\times\mathrm{monon}$ & $2^1 + 2 = 4$ & $2(4)+1 = 9$ & $1{:}2{:}6 = \mathrm{Dimonhexon}$\\
$3{:}1$ & Trio & Trio & $1\times\mathrm{trion}$ & $1^3 + 1 = 2$ & $2(2)+1 = 5$ & $3{:}1{:}6 = \mathrm{Triohexon}$\\
$3{:}2$ & Ditrion & Ditrio & $2\times\mathrm{trion}$ & $2^3 + 2 = 10$ & $2(10)+1 = 21$ & $3{:}2{:}6 = \mathrm{Ditriohexon}$\\
$3{:}3$ & Tritrion & Tritrio & $3\times\mathrm{trion}$ & $3^3 + 3 = 30$ & $2(30)+1 = 61$ & $3{:}3{:}6 = \mathrm{Tritriohexon}$\\
$3{:}4$ & Tetratrion & Tetratrio & $4\times\mathrm{trion}$ & $4^3 + 4 = 68$ & $2(68)+1 = 137$ & $3{:}4{:}6 = \mathrm{Tetratriohexon}$\\
\bottomrule
\end{tabular}
\end{table}

\paragraph{Calibration.}
The count compression basis of \mainref{App.~B.4.1} fixes the ordered triple
\begin{equation}
\mathcal{B}_{\mathrm{curl}} = (3{:}3,\; 3{:}4,\; 1{:}2) = (\mathrm{Tritrio},\; \mathrm{Tetratrio},\; \mathrm{Dimon}),
\label{eq:matter1-count-basis}
\end{equation}
so that any count form $a{:}b{:}c\,{:}{:}\,L$ expands exactly as
\begin{equation}
a{:}b{:}c\,{:}{:}\,L \;\equiv\; \mathrm{Tritrio}^a :: \mathrm{Tetratrio}^b :: \mathrm{Dimon}^c :: \operatorname{out}(L).
\label{eq:matter1-count-expansion}
\end{equation}
In particular, the basic shared hexon is $1{:}1{:}0\,{:}{:}\,6 = \mathrm{Tritrio} :: \mathrm{Tetratrio} :: \mathrm{hexon} \equiv (3{:}3)^1 :: (3{:}4)^1 :: (1{:}2)^0 :: 6$ (\mainref{App.~B.4.2}).

\paragraph{Falsification.}
Invoking $RD_{\AO}(C_f[p{:}q])$ before declaring the certified cycle compression $C_f[p{:}q]_{\AO}$ violates the cycle audit (\mainref{\S5.17.7}). Treating $RD = 2\Pi + 1$ as unconditional when $H_\mu \ne 1$ or when branches are asymmetric violates \mainref{(31)--(32)}.

\aodprovenance{\mainref{(31)--(37)}, \mainref{(40)}, \mainref{(44)--(50)}, \mainref{(168)--(171)}, \mainref{App.~B.1--B.4}.}
\aodliteraturenote{Combinatorial walk support and tesseract encoding: Poyau~\cite{poyau-aod-v4102} (\mainref{\S1.4}, \mainref{App.~H}).}
